\chapter{TEST 1}
\begin{enumerate}
    \item Una empresa compró hace doce meses una maquinaria con un coste de 1.000.000 € con amortización lineal por 4 años, prometiéndose con su devolución de dicha máquina por 200.000 € una vez terminan los 4 años. Al cabo de un año, la empresa decide vender la maquinaria por valor de 900. Calcula el Beneficio Contable.
    \begin{enumerate}
        \item 100.000
        \item 150.000
        \item 200.000
        \item Ninguna de las anteriores
    \end{enumerate}
    \item Una empresa acaba de invertir en 2.000 unidades del producto X, a coste de 150 €/unidad. En el siguiente mes consigue dnegociar mejor con el proveedor y consigue adquirir otras 1.000 unidades del producto X a 135 €/unidad. Si la empresa vende 800 productos a 180 €, ¿cuál será el beneficio de la operación?:
    \begin{enumerate}
        \item 24.000 €
        \item 28.000 €
        \item 36.000 €
        \item Ninguno de los anteriores (porque depende del método de valoración FIFO/LIFO que tenga para sus existencias).
    \end{enumerate}
    \item Una empresa invierte 1.000.000 € en una maquinaria, con una promesa de obtener ingresos de 300.000 € el año 1 y 350.000 € el segundo año. La empresa puede vender la maquinaria al finalizar el segundo año por valor de 500.000 €. Supongamos que la tasa de descuento es dle 9\%. ¿Es rentable invertir en esta maquinaria? ¿por qué?
    \item Elije la opción más lógica
    \begin{enumerate}
        \item El EBITDA mide el resultado bruto de explotación que no se debe confundir con su liquidez.
        \item El EBITDA negativo es algo habitual cuando el proyecto está en sus primeros años de lanzamiento y crecimiento.
        \item Si el EBITDA es negativo, quiere decir que la empresa no es rentable, puesto que el resultado bruto de explotación es negativo e incluso sin hacer frente a intereses e impuestos.
        \item Todas las anteriores
    \end{enumerate}
    \item Respecto a los conceptos financieros ingreso, gasto, cobro y pago, ¿cuál de las siguientes afirmaciones es correcta?
    \begin{enumerate}
        \item Una empresa tiene un ingreso cuando emite una factura a su cliente
        \item Una empresa tiene un ingreso cuando su proveedor paga la factura después de 4 meses de prestar el servicio
        \item Una empresa tiene un gasto cuando paga la factura a su proveedor.
        \item Una empresa tiene un pago cuando recibe una factura de su proveedor sin haberla abonado aún.
    \end{enumerate}
    \item ¿Cuál de las siguientes opciones describe mejor el principio de devengo?
    \begin{enumerate}
        \item Si una empresa presta un servicio en un mes determinado, tanto el ingreso como el cobro se deben contabilizar en ese mismo mes, incluso si el cliente paga la factura meses más tarde.
        \item Si una empresa presta un servicio en un mes determinado, el cobro se debe contabilizar en ese mismo mes, incluso si el cliente paga la factura meses más tarde.
        \item Si una empresa presta un servicio en un mes determinado el ingreso puede ser contabilizado cuando la empresa lo  considere, siempre dentro del año fiscal.
        \item Si una empresa presta un servicio en un mes determinado, el ingreso se debe contabilizar en el mismo mes, incluso si el cliente paga la factura más tarde.
    \end{enumerate}
    \item La empresa Industria S.A. acaba de ser informada por su equipo jurídico en octubre de este año que uno de los litigios judiciales tiene alta probabilidad de pérdida y que la indemnización sería de 50.000 euros, aunque la sentencia no se espera hasta mediados del año próximo 2026. Según el Plan General de Contabilidad, ¿cuál de las siguientes actuaciones es más conforme con los principios contables?
    \begin{enumerate}
        \item Regustar el gasto solo en 2026, cuando se conozca con certeza la pérdida.
        \item Registar el gasto y la provisión correspondiente en 2025, por aplicación tanto del principio de prudencia como principio de devengo.
        \item No registar ninguna provisión, ya que la obligación no es cierta y la sentencia aún no se ha dictado. Lo que hay que registar es un pago con la cuantía notificada por el departamento jurídico.
        \item Registar la provisión correspondiente este año y el gasto el año siguiente
    \end{enumerate}
    \item La empresa Industria S.A. obtiene un préstamo bancario de 250.000 euros a devolver en 4 años, que ingresa íntegramente en su cuenta corriente. ¿Cuál es el efecto inmediato en el balance?
    \begin{enumerate}
        \item Aumenta únicamente el activo en 250.000 €
        \item Aumentan simultáneamente el activo y el pasivo en 250.000 €
        \item Aumentan el activo y el patrimonio neto en 250.000 €
        \item Aumentan simultáneamente el activo y el patrimonio neto en 250.000 €
    \end{enumerate}
    \item Una empresa vende mercancías en diciembre por 40.000 €, pero el cliente no pagará hasta febrero del año siguiente. ¿Qué principio obliga a reconocer el ingreso en diciembre y no en febrero?
    \begin{enumerate}
        \item Devengo
        \item Prudencia
        \item No compensación
        \item Importancia relativa
    \end{enumerate}
    \item Si una empresa obtiene beneficios en el ejercicio, ¿dónde se refleja finalmente ese resultado en el balance?
    \begin{enumerate}
        \item En el activo corriente
        \item En el pasivo corriente
        \item En el patrimonio neto
        \item En los ingresos anticipados
    \end{enumerate}
    \item ¿Cómo se relacionan el balance y la cuenta de resultados?
    \begin{enumerate}
        \item Son documentos totalmente independientes
        \item La cuenta de resultados siempre se elabora antes del balance
        \item El resultado neto de la cuenta de resultados modifica el patrimonio neto del balance
        \item El balance se elabora antes que la cuenta de resultados
    \end{enumerate}
    \item ¿Puede el patrimonio neto ser negativo?
    \begin{enumerate}
        \item Si, cuando la empresa incurre en pérdidas y éstas son mayores que el capital social.
        \item Si, siempre que la empresa incurra en pérdidas
        \item No, nunca
        \item Si, cuando el beneficio neto es superior a las reservas dentro del patrimonio neto.
    \end{enumerate}
    \item ¿Puede el capital social de una empresa aparecer con valor negativo en el balance de situación?
    \begin{enumerate}
        \item Si, si la empresa tiene pérdidas acumuladas superiores al capital inicial
        \item Si, cuando el patrimonio neto es negativo
        \item No, porque el capital social es una cifra legal fijada en los estatutos y no puede ser negativa
        \item No, salvo que la empresa sea una sociedad limitada unipersonal
    \end{enumerate}
    \item Al cierre del ejercicio Industria S.A. obtiene un beneficio de 40.000 € ¿Cómo afecta este resultado al balance de la situación, si aún no se ha distribuido?
    \begin{enumerate}
        \item Aumenta el activo y el pasivo en 40.000 €
        \item Aumenta el patrimonio neto en 40.000 €, sin afectar directamente al activo ni al pasivo
        \item Aumenta el acivo en 40.000 € y reduce el patrimonio neto en el mismo importe
        \item Disminuye el pasivo en 40.000 € y aumenta las reservas.
    \end{enumerate}
    \item Industria S.A. ha presentado en su balance un patrimonio neto inferior al año anterior, pero la cuenta de resultados muestra beneficios (BN). ¿Cuál puede ser la causa?
    \begin{enumerate}
        \item Se han distribuido dividendos mayores que el beneficio obtenido
        \item El activo ha aumentado más que el pasivo
        \item Se ha reducido el capital social mediante ampliación de capital
        \item El capital social es negativo
    \end{enumerate}
    \item Industria S.A. presenta un beneficio de 80.000 € en la cuenta de resultado. En el balance, el activo total aumenta en 200.000 € y el pasivo total aumenta en 120.000 € respecto al ejercicio anterior. ¿Cuál de las siguientes afirmaciones es correcta?
    \begin{enumerate}
        \item El incremento del patrimonio neto es de 80.000 €, correspondiente al beneficio obtenido
        \item El incremento del patrimonio neto es de 200.000 €, ya que el activo ha crecido más que el pasivo
        \item El incremento del patrimonio neto es de 120.000 €, porque el pasivo ha aumentado en esa cantidad.
        \item El patrimonio neto no varía, porque el beneficio solo afecta a la cuenta de resultados.
    \end{enumerate}
    \item Un aumento de los dividendos pagados a los accionistas. ¿Qué efecto produce sobre el beneficio?
    \begin{enumerate}
        \item Ninguna de las otras respuestas
        \item El beneficio no resulta afectado
        \item El beneficio disminuye
        \item El beneficio aumenta
    \end{enumerate}
    \item El registro de los salarios devengados y pendientes de pago, en la fecha de cierre del ejercicio, tiene por objeto:
    \begin{enumerate}
        \item Contabilizar los gastos del ejercicio terminado
        \item Contabilizar las deudas registradas en la fecha de cierre del ejercicio terminado
        \item Reducir el efectivo disponible en la fecha de cierre del ejercicio terminado
        \item Lo indicado en las respuestas a y b
    \end{enumerate}
    \item El pago de 15.000 euros el 15 de diciembre del año 2022, por una campaña de publicidad que se realizará en el primer trimestre del año 2023, afectará al balance de situación al 31 de diciembre de 2022 reduciendo la tesorería en 15.000 euros e incrementando en igual importe:
    \begin{enumerate}
        \item Una cuenta de gastos (627 Gastos de publicidad)
        \item Una cuenta de pasivo (485 Ingresos anticipados)
        \item Una cuenta de ingresos (705 Prestación de servicios)
        \item Una cuenta de activo (480 Gastos anticipados)
        \item Ninguna de las otras respuestas
    \end{enumerate}
    \item Los activos adquiridos por la empresa, en el momento de su incorporación al patrimonio de ésta, se valoran con carácter general por:
    \begin{enumerate}
        \item El coste pagado o comprometido
        \item Cualquier de las otras respuestas
        \item El valor de mercado cuando ésta sea superior al precio pagado por la empresa
        \item La media aritmética del coste pagado y el valor de mercado cuando ésta sea superior al precio pagado por la empresa
    \end{enumerate}
\end{enumerate}